\begin{statement}{A.2.1}{Proved}{The Jacobian fractional ideal for every reduced hypersurface}
Let $\mathcal J\subset\OO_S$ be the Jacobian ideal, generated locally by $F_{x_0},\ldots,F_{x_n}$ modulo $F$. In the total quotient sheaf $\mathcal M_S$ define
\[
 \mathcal J^{-1}=\{q\in\mathcal M_S:q\mathcal J\subset\OO_S\}.
\]
There are canonical identifications
\[
 K\simeq N^*\otimes\mathcal J^{-1},\qquad
 T\simeq N^*\otimes(\mathcal J^{-1}/\OO_S),
\tag{A.2}
\]
under which the first map in \eqref{A:filtration1} is induced by $\OO_S\hookrightarrow\mathcal J^{-1}$. Thus (A.1), with this explicitly defined fractional ideal, is an exact necessary and sufficient criterion for every reduced hypersurface.
\end{statement}
\begin{proof}
The ideal $J$ contains a nonzerodivisor: it is nonzero at every minimal prime, and prime avoidance supplies an element outside all minimal primes. Every $R$-linear map $\varphi:J\to R$ is multiplication by a unique element of the total quotient ring. Indeed, for a nonzerodivisor $a\in J$, the identity $a\varphi(b)=b\varphi(a)$ for $b\in J$ gives the multiplier $q=\varphi(a)/a$. Hence $J^{-1}=\Hom_R(J,R)$ is coherent.

At the generic points, an element $\alpha\in K$ is uniquely $q\,dF$ for a tuple of meromorphic functions $q$. The coefficients of $\alpha$ belong to $R$ exactly when $qF_{x_i}\in R$ for every $i$, which is the condition $q\in J^{-1}$. Conversely such a vector restricts to zero in the cotangent space of every regular component, hence belongs to $K$. The map $[F]\otimes q\mapsto q\,dF|_S$ proves the first identification. Its quotient by $R\,dF$ gives the second. Replacing $F$ by $uF$, with $u$ a holomorphic unit, replaces $dF|_S$ by $u\,dF|_S$ and replaces the conormal generator by $u[F]$. Coordinate changes transform the derivative vector by an invertible Jacobian matrix. These are exactly the transitions in (A.2), so the identifications are canonical.
\end{proof}


% A.2.2 from part_A.tex

